% Vietnamese translation for OI Public Library
% Source-PDF: 动态规划 DYNAMIC PROGRAMMING/动态规划之队列优化.pdf
\documentclass[11pt]{article}
\usepackage{oipl-vn}

\title{Sơ lược về ứng dụng của hàng đợi trong một lớp bài toán đơn điệu}
\author{Tang Ze}
\date{}

\begin{document}
\maketitle

\origtitle{A brief analysis of queues in a class of monotonicity problems}
\sourcepath{Dynamic Programming / Queue Optimization.pdf}

\section{Ví dụ 1: Chọn vị trí xưởng cưa}

Bài CEOI 2004.

\subsection*{Đề bài}

Trên một đường thẳng từ đỉnh núi xuống chân núi có $n$ cây. Chính quyền địa
phương quyết định chặt chúng. Để không lãng phí gỗ, sau khi cây bị chặt, gỗ
phải được vận chuyển tới xưởng cưa. Gỗ chỉ được vận chuyển theo một hướng:
xuống núi. Ở chân núi đã có một xưởng cưa; ngoài ra sẽ xây thêm hai xưởng
cưa mới trên đường núi. Cần quyết định vị trí xây hai xưởng cưa này sao cho
tổng chi phí vận chuyển nhỏ nhất. Giả sử vận chuyển một ki-lô-gam gỗ đi một
mét tốn một xu.

Nhiệm vụ là đọc số cây, trọng lượng và vị trí của chúng, tính chi phí vận
chuyển nhỏ nhất, rồi in kết quả. Dữ liệu có
\[
  2\le n\le 20000,\qquad
  1\le v_i\le 10000,\qquad
  0\le d_i\le 10000.
\]
Cây được đánh số từ đỉnh núi xuống chân núi. Dòng thứ $i+1$ chứa trọng lượng
$v_i$ của cây thứ $i$ và khoảng cách $d_i$ từ cây thứ $i$ tới cây thứ
$i+1$; riêng $d_n$ là khoảng cách từ cây thứ $n$ tới xưởng cưa ở chân núi.
Chi phí đưa toàn bộ gỗ về xưởng ở chân núi nhỏ hơn $2\,000\,000\,000$ xu.

Ví dụ mẫu có đáp án $26$.

\subsection*{Thiết lập}

Trước hết cần nhận ra rằng xây xưởng cưa giữa hai cây kề nhau là vô ích. Nếu
một xưởng nằm giữa hai cây, ta có thể dời nó lên đúng vị trí cây gần nhất ở
phía trên; khi đó chi phí vận chuyển gỗ từ phía trên giảm, còn chi phí của
phía dưới không đổi. Vì vậy, chỉ cần xét các vị trí tại cây.

Để tiện thảo luận, giả sử ở chân núi cũng có một cây giả, đánh số $n+1$, với
$v[n+1]=d[n+1]=0$. Định nghĩa
\[
  \operatorname{sumw}[i]=\sum_{j=1}^{i} w[j],
  \qquad
  \operatorname{sumd}[i]=\sum_{j=1}^{i-1} d[j],
\]
trong đó $\operatorname{sumd}[1]=0$.

Gọi $c[i]$ là chi phí nếu xây một xưởng cưa tại cây thứ $i$ và vận chuyển
toàn bộ cây từ $1$ đến $i$ tới đó. Khi đó
\[
  c[1]=0,\qquad
  c[i]=c[i-1]+\operatorname{sumw}[i-1]d[i-1].
\]
Gọi $w[j,i]$ là chi phí nếu xây xưởng tại cây thứ $i$ và vận chuyển các cây
từ $j$ đến $i$ tới đó. Ta có
\[
  w[j,i]
  =
  c[i]-c[j-1]
  -\operatorname{sumw}[j-1]\bigl(\operatorname{sumd}[i]
    -\operatorname{sumd}[j-1]\bigr),
\]
và $w[j,i]=0$ khi $i\le j$.

Tất cả các mảng $\operatorname{sumw}$, $\operatorname{sumd}$ và $c$ được tính
trong $O(n)$; sau đó mỗi giá trị $w[j,i]$ được tính trong $O(1)$.

Gọi $f[i]$ là chi phí nhỏ nhất khi xưởng cưa thứ hai được xây ở cây thứ $i$.
Khi đó
\[
  f[i]=
  \min_{1\le j\le i-1}
  \{c[j]+w[j+1,i]+w[i+1,n+1]\}.
\]
Thuật toán trực tiếp mất $O(n^2)$, không thể dùng cho kích thước dữ liệu này.

\subsection*{Tính đơn điệu của quyết định}

Ta chứng minh nhận xét sau. Nếu khi xây xưởng cưa thứ hai ở vị trí $i$, vị
trí tối ưu của xưởng thứ nhất là $j$, thì khi xây xưởng thứ hai ở vị trí
$i+1$, vị trí tối ưu của xưởng thứ nhất không nhỏ hơn $j$.

Giả sử $j$ là vị trí nhỏ nhất để $f[i]$ đạt tối ưu:
\[
  f[i]=c[j]+w[j+1,i]+w[i+1,n+1].
\]
Với mọi $k<j$, ta có
\[
  c[k]+w[k+1,i]+w[i+1,n+1]
  >
  c[j]+w[j+1,i]+w[i+1,n+1],
\]
nên
\[
  c[k]+w[k+1,i]>c[j]+w[j+1,i].
\]
Khi mở rộng từ $i$ sang $i+1$, phần gỗ từ $k+1$ đến $i$ nặng hơn phần gỗ từ
$j+1$ đến $i$ vì $k<j$. Do đó sau khi cộng thêm chi phí vận chuyển qua đoạn
$d[i]$, quyết định $k$ vẫn kém $j$:
\[
  c[k]+w[k+1,i+1]
  >
  c[j]+w[j+1,i+1].
\]
Cộng thêm cùng một lượng $w[i+2,n+1]$, ta thấy khi tính $f[i+1]$, quyết định
$k$ vẫn kém quyết định $j$. Vì mọi $k<j$ đều bị loại, quyết định tối ưu đầu
tiên của $f[i+1]$ phải lớn hơn hoặc bằng $j$.

\subsection*{Cải tiến bằng hàng đợi}

Đặt $s[k,i]$ là giá trị của hàm mục tiêu khi biến quyết định bằng $k$:
\[
  s[k,i]=c[k]+w[k+1,i]+w[i+1,n+1].
\]
Với $k_1<k_2$, ta có
\[
\begin{aligned}
  s[k_1,i]-s[k_2,i]
  &=
    \operatorname{sumd}[i]\bigl(\operatorname{sumw}[k_2]
      -\operatorname{sumw}[k_1]\bigr) \\
  &\quad-
    \bigl(\operatorname{sumw}[k_2]\operatorname{sumd}[k_2]
      -\operatorname{sumw}[k_1]\operatorname{sumd}[k_1]\bigr).
\end{aligned}
\]
Nếu $s[k_1,i]-s[k_2,i]<0$, thì
\[
  \frac{
    \operatorname{sumw}[k_2]\operatorname{sumd}[k_2]
    -\operatorname{sumw}[k_1]\operatorname{sumd}[k_1]}
  {\operatorname{sumw}[k_2]-\operatorname{sumw}[k_1]}
  >
  \operatorname{sumd}[i].
\]
Ký hiệu vế trái là $g[k_1,k_2]$. Khi
$g[k_1,k_2]>\operatorname{sumd}[i]$, quyết định $k_1$ tốt hơn $k_2$ tại
trạng thái $i$.

Do quyết định tối ưu đơn điệu, ta duy trì một hàng đợi các quyết định
$k_1<k_2<\cdots<k_p$ sao cho
\[
  g[k_1,k_2]<g[k_2,k_3]<\cdots<g[k_{p-1},k_p].
\]
Trước khi tính $f[i]$, nếu
$g[k_1,k_2]<\operatorname{sumd}[i]$, nghĩa là $k_1$ không tốt bằng $k_2$,
ta xóa $k_1$ khỏi đầu hàng đợi, và lặp lại cho đến khi đầu hàng đợi là quyết
định tốt nhất. Khi đó
\[
  f[i]=c[k_1]+w[k_1+1,i]+w[i+1,n+1]
\]
được tính trong $O(1)$.

Sau khi tính $f[i]$, ta chèn quyết định mới $i$ vào cuối hàng đợi. Nếu
\[
  g[k_{p-1},k_p]>g[k_p,i],
\]
thì $k_p$ trở nên vô ích: trước khi $k_p$ tốt hơn $k_{p-1}$, quyết định $i$
đã tốt hơn $k_p$. Vì vậy ta xóa $k_p$ và lặp lại cho đến khi thứ tự các giá
trị $g$ được khôi phục. Mỗi quyết định vào hàng đợi một lần và ra hàng đợi
một lần, nên tổng chi phí duy trì hàng đợi là $O(n)$. Do đó thuật toán chạy
trong $O(n)$.

Điều kiện quan trọng trong ví dụ này là trọng lượng mỗi cây đều dương. Nhờ
đó mảng tiền tố trọng lượng đơn điệu, và lập luận ``đoạn từ $j+1$ đến $i$
nhẹ hơn đoạn từ $k+1$ đến $i$ khi $k<j$'' mới đúng.

\section{Ví dụ 2: Xây kho}

Bài tuyển chọn tỉnh Chiết Giang năm 2007.

\subsection*{Đề bài}

Công ty L có $N$ nhà máy nằm trên một ngọn núi, đánh số từ cao xuống thấp:
nhà máy $1$ ở đỉnh núi, nhà máy $N$ ở chân núi. Công ty thường để sản phẩm
ngoài trời để tiết kiệm chi phí, nhưng do sắp có mưa lớn, họ cần xây một số
kho tại các nhà máy. Nhà máy thứ $i$ hiện có $P_i$ sản phẩm; chi phí xây kho
tại đó là $C_i$. Sản phẩm của nhà máy không có kho phải được vận chuyển
xuống các kho ở phía dưới, tức chỉ có thể đi tới nhà máy có chỉ số lớn hơn.
Chi phí vận chuyển một sản phẩm qua một đơn vị khoảng cách là $1$, và mỗi kho
có sức chứa đủ lớn.

Dữ liệu gồm khoảng cách $X_i$ từ nhà máy $1$ đến nhà máy $i$ ($X_1=0$), số
sản phẩm $P_i$ và chi phí xây kho $C_i$. Cần tìm phương án xây kho sao cho
tổng chi phí xây dựng và vận chuyển nhỏ nhất. Kích thước dữ liệu có thể tới
$N\le 10^6$, các số nằm trong kiểu số nguyên có dấu $32$ bit, và kết quả
trung gian không vượt quá kiểu số nguyên có dấu $64$ bit.

\subsection*{Phân tích ban đầu}

Đặt
\[
  \operatorname{sump}[i]=\sum_{j=1}^{i}P_j,
\]
và $\operatorname{sumw}[i]$ là chi phí vận chuyển toàn bộ sản phẩm từ nhà máy
$1$ đến $i$ về nhà máy $i$. Gọi $w[i,j]$ là chi phí vận chuyển sản phẩm từ
nhà máy $i$ đến $j$ về nhà máy $j$. Khi đó
\[
  w[i,j]=\operatorname{sumw}[j]-\operatorname{sumw}[i-1]
  -\operatorname{sump}[i-1]\bigl(x[j]-x[i-1]\bigr).
\]
Các giá trị tiền tố được tính trong $O(n)$, và mỗi $w[i,j]$ được tính trong
$O(1)$.

Gọi $f[i]$ là chi phí nhỏ nhất khi xây một kho tại nhà máy $i$ và xử lý xong
các nhà máy từ $1$ đến $i$. Đáp án là $f[n]$, với
\[
  f[i]=\min_{0\le k\le i-1}
    \{f[k]+w[k+1,i]+c[i]\},
  \qquad f[0]=0.
\]
Thuật toán trực tiếp có độ phức tạp $O(n^2)$. Dưới đây dùng cùng tư tưởng như
ví dụ thứ nhất để giảm xuống $O(n)$.

\subsection*{Đơn điệu và tối ưu hóa}

Giả sử quyết định tối ưu khi tính $f[i]$ là $k$. Với mọi $j<k$, từ tính tối
ưu của $k$ ta có
\[
  f[k]+w[k+1,i]+c[i] < f[j]+w[j+1,i]+c[i].
\]
Viết lại bằng các mảng tiền tố:
\[
\begin{aligned}
  f[k]-\operatorname{sumw}[k]
  +\operatorname{sump}[k](x[i]-x[k])
  &<
  f[j]-\operatorname{sumw}[j]
  +\operatorname{sump}[j](x[i]-x[j]).
\end{aligned}
\]
Vì $\operatorname{sump}[k]\ge\operatorname{sump}[j]$ và $x[i+1]\ge x[i]$,
bất đẳng thức vẫn đúng khi thay $x[i]$ bằng $x[i+1]$. Do đó quyết định $k$
cũng tốt hơn $j$ khi xét trạng thái $i+1$. Suy ra quyết định tối ưu đơn điệu
theo $i$.

Đặt
\[
  s[k,i]=f[k]+w[k+1,i]+c[i].
\]
Với $k_1<k_2$,
\[
\begin{aligned}
  s[k_1,i]-s[k_2,i]
  &=
  \bigl(f[k_1]+\operatorname{sumw}[k_1]
    +\operatorname{sump}[k_1]x[k_1]\bigr)\\
  &\quad-
  \bigl(f[k_2]+\operatorname{sumw}[k_2]
    +\operatorname{sump}[k_2]x[k_2]\bigr)\\
  &\quad-
  x[i]\bigl(\operatorname{sump}[k_1]-\operatorname{sump}[k_2]\bigr).
\end{aligned}
\]
Khi $s[k_1,i]-s[k_2,i]<0$, ta nhận được một ngưỡng chỉ phụ thuộc vào
$k_1,k_2$:
\[
  g[k_1,k_2]=
  \frac{
  f[k_1]+\operatorname{sumw}[k_1]+\operatorname{sump}[k_1]x[k_1]
  -
  f[k_2]-\operatorname{sumw}[k_2]-\operatorname{sump}[k_2]x[k_2]}
  {\operatorname{sump}[k_1]-\operatorname{sump}[k_2]}
  >x[i].
\]
Vì vậy có thể dùng hàng đợi đơn điệu tương tự ví dụ trước: xóa đầu hàng đợi
khi quyết định kế tiếp đã tốt hơn tại $x[i]$, dùng đầu hàng đợi để tính
$f[i]$, rồi chèn quyết định $i$ ở cuối và xóa các quyết định bị nó làm cho
vô dụng. Tổng thời gian vẫn là $O(n)$.

\section{Ví dụ 3: Vũ điệu lộng lẫy}

Bài NOI 2005 ngày 1.

\subsection*{Đề bài}

Cho một ma trận $N$ hàng, $M$ cột; một số ô là chướng ngại vật. Một người bắt
đầu trượt từ một ô cho trước. Mỗi lần trượt có một hướng và được đi liên tục
tối đa $c$ ô theo hướng đó, cũng có thể đứng yên. Các lần trượt có thể có
giá trị $c$ khác nhau. Trong khi trượt không được chạm chướng ngại vật. Cho
theo thứ tự $K$ lần trượt, mỗi lần gồm hướng đông, nam, tây hoặc bắc và giá
trị $c$, hãy tìm quãng đường dài nhất có thể trượt được, tính bằng số ô.

Dữ liệu có
\[
  1\le N,M\le 200,\qquad K\le 200,\qquad
  \sum_{i=1}^{K}c_i\le 40000.
\]

\subsection*{Quy hoạch động trực tiếp}

Gọi $f(k,x,y)$ là quãng đường dài nhất sau $k$ lần trượt và đang ở ô
$(x,y)$. Chỉ xét hướng đông; ba hướng còn lại suy ra tương tự. Ta có
\[
  f(0,\operatorname{startx},\operatorname{starty})=0,
\]
và
\[
\begin{aligned}
  f(k,x,y)=\max\{&
    f(k-1,x,y),\\
    &f(k-1,x,y-1)+1,\\
    &f(k-1,x,y-2)+2,\ldots,\\
    &f(k-1,x,y')+y-y'\},
\end{aligned}
\]
trong đó $y'$ là vị trí xa nhất bên trái còn có thể trượt tới mà không vượt
quá $c_k$ ô và không đi qua chướng ngại vật.

Số trạng thái là $O(KMN)$, nhưng mỗi lần chuyển có thể mất
$O(\max\{M,N\})$, nên tổng thời gian là
$O(KMN\max\{M,N\})$, quá lớn.

\subsection*{Biến đổi công thức}

Không xét chướng ngại vật trước. Khi tính $f(k,x,y)$ và $f(k,x,y+1)$, rất
nhiều trạng thái cũ bị xét lặp lại. Với một hàng cố định và hướng đông, công
thức có dạng
\[
  f(k,x,y)=
  \max\{f(k-1,x,t)+y-t\mid y-c_k\le t\le y\}.
\]
Đặt
\[
  a_t=f(k-1,x,t)-t+1.
\]
Khi đó các ứng viên chỉ khác nhau bởi cùng một lượng phụ thuộc vào $y$, nên
bài toán trở thành lấy giá trị lớn nhất của một đoạn liên tiếp trên dãy
$a_t$.

Khi thêm chướng ngại vật, mỗi hàng được chia thành các đoạn không có vật cản.
Trong quá trình quét từ trái sang phải, tập chỉ số hợp lệ luôn là một đoạn:
mỗi bước thêm một giá trị ở cuối, đôi khi xóa vài giá trị ở đầu vì độ dài
vượt quá $c_k+1$, và khi gặp vật cản thì xóa sạch đoạn hiện tại. Có thể dùng
cây đoạn để hỏi giá trị lớn nhất, đạt thời gian
$O(KMN\log\max\{M,N\})$, nhưng cách này vừa dài vừa có hằng số lớn.

\subsection*{Hàng đợi đơn điệu}

Tính chất của đoạn đang quét phù hợp với hàng đợi: chỉ thêm ở một đầu và xóa
ở đầu kia. Vấn đề là cần lấy giá trị lớn nhất trong hàng đợi. Ta lưu hàng đợi
theo dạng đơn điệu giảm của giá trị $a$.

Nếu trong hàng đợi đã có hai phần tử $a_2$ và $a_3$ với $a_2\le a_3$, thì
$a_2$ không bao giờ cần dùng. Khi $a_3$ đã xuất hiện, trước khi $a_2$ bị xóa
thì $a_3$ cũng vẫn còn trong hàng đợi; vì $a_2\le a_3$, giá trị lớn nhất
không thể là $a_2$. Vì vậy khi chèn một phần tử mới $a_i$, ta xóa từ cuối
hàng đợi mọi phần tử không lớn hơn $a_i$, rồi mới chèn $i$.

Khi xử lý ô hiện tại, nếu chỉ số ở đầu hàng đợi cách ô hiện tại quá $c_k$,
ta xóa nó khỏi đầu. Giá trị lớn nhất trong cửa sổ hiện tại luôn nằm ở đầu
hàng đợi, nên truy vấn mất $O(1)$. Mỗi phần tử được chèn nhiều nhất một lần
và xóa nhiều nhất một lần trên mỗi hàng hoặc cột, do đó tổng thời gian cho
một lần trượt là $O(MN)$, và toàn bộ thuật toán chạy trong $O(KMN)$.

Với hướng đông, mã giả có thể viết như sau:
\[
\begin{array}{l}
\text{khởi tạo } f[0][x][y] = -\infty \text{ cho mọi ô,}\\
f[0][\operatorname{startx}][\operatorname{starty}]=0;\\
\text{for } k=1\ldots K:\\
\quad \text{for } x=1\ldots n:\\
\quad\quad \text{xóa rỗng hàng đợi;}\\
\quad\quad \text{for } y=1\ldots m:\\
\quad\quad\quad \text{nếu }(x,y)\text{ là chướng ngại, đặt }
  f[k][x][y]=-\infty\text{ và xóa hàng đợi;}\\
\quad\quad\quad \text{nếu đầu hàng đợi }<y-c_k,\text{ xóa đầu;}\\
\quad\quad\quad \text{xóa cuối khi giá trị ở cuối không lớn hơn }
  f[k-1][x][y]-y+1;\\
\quad\quad\quad \text{chèn }y;\\
\quad\quad\quad f[k][x][y]=
  f[k-1][x][\operatorname{front}]+y-\operatorname{front}.\\
\end{array}
\]
Các hướng còn lại chỉ thay đổi thứ tự quét và cách tính chỉ số.

\section{Kết luận}

Nhìn lại toàn bộ quá trình, ta thường bắt đầu bằng một lời giải quy hoạch
động tự nhiên. Khi mỗi chuyển trạng thái quá đắt, bước quan trọng là nhận ra
các phép tính lặp và tìm cấu trúc đơn điệu nằm trong biến quyết định hoặc
trong cửa sổ chuyển trạng thái. Cây đoạn là một công cụ tổng quát tốt, nhưng
trong nhiều trường hợp hàng đợi đơn điệu đơn giản hơn, nhanh hơn và ít lỗi
hơn khi cài đặt.

Điều này cho thấy trong xu hướng mới của các kỳ thi tin học, vai trò của cấu
trúc dữ liệu cơ bản không hề giảm đi; ngược lại, chúng ngày càng quan trọng.
Khi đã tìm được một thuật toán tốt nhưng việc cài đặt quá phức tạp, nên thử
đổi góc nhìn và dùng các cấu trúc dữ liệu cơ bản. Nhiều khi cách này đem lại
hiệu quả gấp bội.

\end{document}
